%% toc: P1 8 8.3,8.4,8.6,8.7
\section{Proofs for the equivalence glyphs and the colon-equals sign}
\label{sec:annex-p01-w02b-p01-013-equivalence-audit}

\begin{annexitem}{annex:p01-w02b-p01-013-congruence-definable}{thm:p01-w02b-p01-013-congruence-definable}{Congruence predicates are definitional}
\begin{proof}
Fix \(m\geq1\) and \(0\leq r<m\).

\emph{Values of the auxiliary terms.}  Under every assignment the term
\(\underline r\) has the value \(r\): for \(r=0\) it is the constant \(0\);
for \(r\geq1\) it is a sum of \(r\) copies of the constant \(1\), whose value
is \(r\) by induction on the number of summands.  If the assignment sends
\(y\) to \(k\), the term \(m\cdot y\) has the value \(mk\), again by induction
on the number \(m\) of summands \(y\).  So under an assignment sending \(x\)
to \(n\) and \(y\) to \(k\), the atom \(x=\underline r+m\cdot y\) is true in
\(\mathbb N\) iff \(n=r+mk\).

\emph{Part~1.}  By the satisfaction clause for \(\exists\), the formula
\(\exists y\,(x=\underline r+m\cdot y)\) is true under \(x\mapsto n\) iff
\(n=r+mk\) for some \(k\in\mathbb N\).  If \(n=r+mk\), then \(n-r=mk\) and
\(m\) divides \(n-r\).  Conversely, let \(n-r=mq\) for an integer \(q\).
Since \(n\geq0\) and \(r<m\), we have \(mq=n-r>-m\), so \(q>-1\), and \(q\)
is an integer, hence \(q\geq0\).  Then \(k=q\in\mathbb N\) and \(n=r+mk\).

\emph{Part~2.}  By the division algorithm in \(\mathbb N\) there are
\(q,\rho\in\mathbb N\) with \(n=qm+\rho\) and \(\rho<m\); then \(m\) divides
\(n-\rho\), so the condition holds for \(r=\rho\).  If it holds for
\(r_1,r_2<m\), then \(m\) divides both \(n-r_1\) and \(n-r_2\), hence divides
their difference \(r_2-r_1\).  Since \(|r_2-r_1|<m\), this difference is
\(0\), so \(r_1=r_2\).

\emph{Part~3.}  Let \(y\) be a variable that does not occur in \(\chi\).  The
proof is by induction on \(\chi\) and shows the claim for every subformula
\(\chi'\) of \(\chi\) and every assignment \(s\); the translation of a
subformula uses the same \(y\), which occurs in no subformula either.

If \(\chi'\) is an atom of the original language, then
\(\chi'^{\circ}=\chi'\), and \(\mathbb N^{+}\) and \(\mathbb N\) interpret
its symbols identically, so the two satisfaction conditions coincide.

If \(\chi'\) is \(C_{m',r'}(t)\), let \(n\) be the value of \(t\) under
\(s\).  By the interpretation in \(\mathbb N^{+}\), \(\mathbb
N^{+}\entails\chi'[s]\) iff \(m'\) divides \(n-r'\).  The translation is
\(\exists y\,(t=\underline{r'}+m'\cdot y)\).  For every \(k\), the
assignment \(s[y\mapsto k]\) gives \(t\) the value \(n\), because \(y\) does
not occur in \(t\).  The computation of the first paragraph, with \(t\) in
place of \(x\), therefore shows that \(\mathbb N\entails\chi'^{\circ}[s]\)
iff \(n=r'+m'k\) for some \(k\in\mathbb N\), and part~1 (for \(m',r'\))
turns this into the condition that \(m'\) divides \(n-r'\).

If \(\chi'\) is \(\neg\chi_1\) or \(\chi_1\circ\chi_2\) for a binary
connective \(\circ\), then \(\chi'^{\circ}\) is \(\neg\chi_1^{\circ}\) or
\(\chi_1^{\circ}\circ\chi_2^{\circ}\).  The truth clause for the connective
in both structures uses only the truth values of the immediate subformulas
under \(s\), and these agree by the induction hypothesis.

If \(\chi'\) is \(Qz\,\chi_1\) with \(Q\in\{\forall,\exists\}\), then
\(\chi'^{\circ}\) is \(Qz\,\chi_1^{\circ}\); here \(z\neq y\) because \(y\)
does not occur in \(\chi\).  Both structures have the domain \(\mathbb N\).
For each \(n\in\mathbb N\) the induction hypothesis for \(\chi_1\) under
\(s[z\mapsto n]\) gives \(\mathbb N^{+}\entails\chi_1[s[z\mapsto n]]\) iff
\(\mathbb N\entails\chi_1^{\circ}[s[z\mapsto n]]\).  The quantifier clause
asks for this for every \(n\) (for \(\forall\)) or some \(n\) (for
\(\exists\)) in both structures, so the two sides agree.

Taking \(\chi'=\chi\) gives part~3.
\end{proof}
\end{annexitem}

\begin{annexitem}{annex:p01-w02b-p01-013-identically-equal}{thm:p01-w02b-p01-013-identically-equal}{Identical equality is pointwise}
\begin{proof}
\emph{Part~1.}  Two functions with domain \(D\) and the same codomain are
equal iff they take the same value at every point of \(D\).  The function
\(\kappa_c\) takes the value \(c\) at every point, so \(h=\kappa_c\) iff
\(h(y)=c\) for every \(y\in D\), which is the definition of \(h\equiv c\) in
\cref{def:p01-w02b-p01-013-identically-equal}.

\emph{Part~2.}  The set of points at which \(h'\) and \(h\) differ is
contained in \(\{y_0\}\), since they agree elsewhere by hypothesis.  It
equals \(\{y_0\}\): \(h(y_0)=c\) because \(h\equiv c\), and
\(h'(y_0)=c'\neq c\).  This set belongs to \(\mathcal D\) and has measure
\(0\), so \(h'=h\) \(\mu\)-almost everywhere.  Finally \(h'(y_0)=c'\neq c\)
is a point of \(D\) at which \(h'\) does not take the value \(c\), so
\(h'\not\equiv c\) by part~1.
\end{proof}
\end{annexitem}

\begin{annexitem}{annex:p01-w02b-p01-013-contradictory}{thm:p01-w02b-p01-013-contradictory}{Contradictory pairing and nonequivalence}
\begin{proof}
Write \(P(c)\) and \(Q(c)\) for the truth values of \(P\) and \(Q\) in the
case \(c\in C\), and \(\neg P\) for the proposition whose value in each case
is the other value from \(P(c)\).

\emph{Part~1.}  Suppose \(P\) and \(Q\) are contradictories over \(C\).
Since \(C\) is nonempty, choose \(c\in C\).  Exactly one of \(P(c)\) and
\(Q(c)\) is true, so \(P(c)\neq Q(c)\).  Hence \(P\) and \(Q\) do not have
the same value in every case, and they are not equivalent over \(C\).

\emph{Part~2.}  Let \(c_0\in C\) be a case with \(P(c_0)\) and \(Q(c_0)\)
both false.  In \(c_0\) neither proposition is true, so it is false that
exactly one of them is true in every case; \(P\) and \(Q\) are not
contradictories over \(C\).  If moreover \(c_1\in C\) is a case in which
exactly one of them is true, then \(P(c_1)\neq Q(c_1)\), so \(P\) and \(Q\)
are not equivalent over \(C\) either.

\emph{Part~3.}  For a single case \(c\), exactly one of \(P(c)\) and \(Q(c)\)
is true iff \(Q(c)\neq P(c)\), since there are only two truth values; and
\(Q(c)\neq P(c)\) iff \(Q(c)=(\neg P)(c)\) by the definition of \(\neg P\).
Taking this equivalence in every case \(c\in C\): \(Q\) is a contradictory
of \(P\) over \(C\) iff \(Q(c)=(\neg P)(c)\) for every \(c\in C\), which is
equivalence of \(Q\) and \(\neg P\) over \(C\).
\end{proof}
\end{annexitem}

\begin{annexitem}{annex:p01-w02b-p01-013-colon-equals}{thm:p01-w02b-p01-013-colon-equals}{Definition, update and assignment differ}
\begin{proof}
\emph{Part~1.}  First, for every term \(u\) of \(\mathcal L\cup\{c\}\) and
every assignment \(s\), the value of \(u^{t}\) in \(\mathcal M\) under \(s\)
equals the value of \(u\) in \(\mathcal M^{c}\) under \(s\).  This is shown
by induction on \(u\).  If \(u\) is a variable or a constant of
\(\mathcal L\), then \(u^{t}=u\), and \(\mathcal M\) and \(\mathcal M^{c}\)
interpret it alike.  If \(u\) is \(c\), then \(u^{t}=t\); the value of \(t\)
in \(\mathcal M\) is \(t^{\mathcal M}\), since \(t\) is closed, and the value
of \(c\) in \(\mathcal M^{c}\) is \(t^{\mathcal M}\) by the definition of
\(\mathcal M^{c}\).  If \(u\) is \(g(u_1,\ldots,u_k)\) for a function symbol
\(g\) of \(\mathcal L\), then \(u^{t}=g(u_1^{t},\ldots,u_k^{t})\), the two
structures interpret \(g\) alike, and the arguments have equal values by the
induction hypothesis.

Next, by induction on \(\chi\), for every assignment \(s\),
\(\mathcal M^{c}\entails\chi[s]\) iff \(\mathcal M\entails\chi^{t}[s]\).  If
\(\chi\) is an atom \(R(u_1,\ldots,u_k)\) or \(u_1=u_2\), then \(\chi^{t}\)
is the same atom with each \(u_i\) replaced by \(u_i^{t}\); the relation
symbol and identity are interpreted alike in both structures, and the
argument values agree by the term step.  If \(\chi\) is \(\neg\chi_1\) or
\(\chi_1\circ\chi_2\) for a binary connective \(\circ\), then \(\chi^{t}\) is
\(\neg\chi_1^{t}\) or \(\chi_1^{t}\circ\chi_2^{t}\), and the truth clause
uses only the truth values of the immediate subformulas under \(s\), which
agree by the induction hypothesis.  If \(\chi\) is \(Qz\,\chi_1\) with
\(Q\in\{\forall,\exists\}\), then \(\chi^{t}\) is \(Qz\,\chi_1^{t}\): since
\(t\) is closed, the replacement puts no variable into the scope of \(Qz\).
Both structures have the same domain, and for each element \(d\) of it the
induction hypothesis under \(s[z\mapsto d]\) gives
\(\mathcal M^{c}\entails\chi_1[s[z\mapsto d]]\) iff
\(\mathcal M\entails\chi_1^{t}[s[z\mapsto d]]\).  The quantifier clause asks
for this for every \(d\) or for some \(d\) in both structures, so the two
sides agree.

\emph{Part~2.}  Two functions on \(D\) are equal iff they agree at every
point of \(D\), so each identity is checked pointwise.

For \(f[a:=f(a)]=f\): at \(a\) the left side has the value \(f(a)\); at
\(y\neq a\) it has the value \(f(y)\).

For \(f[a:=b][a:=b']=f[a:=b']\): at \(a\) both sides have the value \(b'\),
the outer update on the left being made at \(a\).  At \(y\neq a\) the left
side has the value of \(f[a:=b]\) at \(y\), which is \(f(y)\), and so does
the right side.

For \(a\neq a'\) and \(f[a:=b][a':=b']=f[a':=b'][a:=b]\): at \(a'\) the left
side has the value \(b'\) from its outer update, and the right side has the
value of \(f[a':=b']\) at \(a'\), which is \(b'\), because its outer update is
at \(a\neq a'\).  At \(a\) the left side has the value of \(f[a:=b]\) at
\(a\), which is \(b\), because its outer update is at \(a'\neq a\); the right
side has the value \(b\) from its outer update.  At \(y\notin\{a,a'\}\) both
sides have the value \(f(y)\).

For the last clause: if \(f(a)=b\), then \(f[a:=b]=f[a:=f(a)]=f\) by the
first identity.  If \(f[a:=b]=f\), then the two functions agree at \(a\),
so \(b=f[a:=b](a)=f(a)\).

\emph{Part~3.}  Let \(\sigma\) be a state and let
\(\sigma'=\sigma[x:=\sigma(x)+1]\).  Then
\(\sigma'(x)=\sigma(x)+1\neq\sigma(x)\), since no natural number equals its
successor, so \(\sigma'\neq\sigma\).  The fact used here is proved as
follows: if \(n=n+1\) for some \(n\in\mathbb N\), then subtracting \(n\)
from both sides gives \(0=1\), which is false.  This also proves the second
claim of part~3.
\end{proof}
\end{annexitem}
